\section{Conclusion and Future Improvements}
\label{sec:conclusion_future_improvements}
\subsection{Conclusion}

This project successfully demonstrated the design and development of a Line Following Robot capable of autonomous navigation and obstacle detection. By integrating infrared sensors, ultrasonic sensors, an Arduino Uno microcontroller, and a motor control system, the robot was able to follow a predefined path while responding to environmental conditions.

A systems engineering approach was applied throughout the project to support requirement identification, system modelling, implementation, and testing. The use of SysML diagrams provided a clear representation of both the structural and behavioural aspects of the system, helping to improve understanding of component interactions and overall system functionality.

The testing results showed that the robot was able to perform the intended line-following and obstacle-detection tasks under normal operating conditions. Therefore, the project achieved its primary objective of developing a functional robotic prototype while applying key systems engineering principles.

\subsection{Future Improvements}

Although the prototype met the project requirements, several improvements could be implemented in future developments. One possible enhancement is the integration of a PID control algorithm to improve movement accuracy and provide smoother line-following performance. Additional infrared sensors could also be incorporated to increase detection reliability, particularly when navigating sharp curves or complex paths.

The mechanical design could be further optimized by improving chassis stability, wheel alignment, and component placement. These modifications may contribute to better overall performance and durability of the robot.

Furthermore, wireless communication technologies such as Bluetooth or Wi-Fi could be added to enable remote monitoring and control. Future versions may also incorporate camera-based vision systems or machine learning techniques to support more advanced navigation and obstacle-avoidance capabilities.

Overall, the proposed improvements would increase the robot’s accuracy, adaptability, and functionality, making it more suitable for complex autonomous applications.
